%\documentclass[10pt,dvipsnames,svgnames]{beamer}
\documentclass[aspectratio=169, lualatex,  10pt,dvipsnames,svgnames]{beamer} % handout,
\makeatletter\def\input@path{{theme/}}\makeatother\usetheme{cipher}

\title{Computer Security: Introduction}
\author{Markulf Kohlweiss}
\subject{Introduction to Applied Cryptography covering cryptographic primitives, protocols, security goals, and real-world applications from basic concepts to advanced topics like TLS and zero-knowledge proofs.}
\keywords{cryptography, security, encryption, authentication, hash functions, TLS, protocols, provable security}
\institute{University of Edinburgh, School of Informatics}
\instituteimage{images/edu-white.png}
\date{\today}
\coversubtitle{INFR10067\\Fall 2025}
\coverpartname{Cryptography}
\covertopicname{Introduction to Protocols}
\coverwebsite{}

\iffalse
\mode<presentation>
{
  \usetheme{default}
  \usecolortheme{default}
  \usefonttheme{default}
  \setbeamertemplate{navigation symbols}{}
  \setbeamertemplate{caption}[empty]
  \setbeamertemplate{footline}[frame number]
}

\setbeamertemplate{blocks}[rounded][shadow=true]
\setbeamercolor{block title}{fg=black, bg=Grey!50}
\setbeamercolor{block body}{fg=black, bg=Grey!25}
\setbeamercolor{block title alerted}{fg=black, bg=Red!20}
\setbeamercolor{block body alerted}{fg=black, bg=Red!15}
\setbeamercolor{block title example}{fg=black, bg=SeaGreen!20}
\setbeamercolor{block body example}{fg=black, bg=Green!15}
\fi

\usepackage[english]{babel}
\usepackage[utf8x]{inputenc}

\usepackage[lambda,advantage,operators,sets,adversary,landau,probability,notions,logic,ff,mm,primitives,events,complexity,asymptotics,keys]{cryptocode}

\usepackage{graphicx}
\usepackage{tikz}
\usetikzlibrary{arrows.meta}
\usetikzlibrary{shapes.geometric}
\pgfdeclarelayer{background}
\pgfsetlayers{background,main}

\usepackage{scrextend} 

\usepackage{fancybox}

\usepackage{pgfpages}

\usepackage{caption}
\usepackage{subcaption}

\usepackage{mdframed}
\usepackage{xspace}
\newcommand\Fontvi{\fontsize{5}{5.2}\selectfont}
\definecolor{ao}{rgb}{0.0, 0.6, 0.0}
\definecolor{bleu}{rgb}{0.2,0.2,0.7}

\newcommand{\Inred}[1]{\textcolor{BrickRed}{#1}}
\newcommand{\Inblue}[1]{\textcolor{bleu}{#1}}
\newcommand{\Ingreen}[1]{\textcolor{SeaGreen}{#1}}

\newcommand{\V}[1]{\texttt{V\textsubscript{{#1}}}\xspace}
\newcommand{\Vs}[1]{$\texttt{V}'_\texttt{{#1}}$\xspace}
\newcommand{\BB}{\texttt{BB}\xspace}
\newcommand{\EA}{\texttt{EA}\xspace}
\newcommand{\A}{\adv\xspace}
\newcommand{\F}[1]{$\mathcal{F}_\textsc{{#1}}$\xspace}
\renewcommand{\S}{\sdv\xspace}
\newcommand{\E}{$\mathcal{E}$\xspace}
\renewcommand{\P}{$P$\xspace}
\newcommand{\D}{\ddv\xspace}

\newcommand{\Pt}[2]{$\Pi^{{#2}}_{\textsc{{#1}}}$\xspace}

\newcommand{\Enc}{\texttt{Enc}}
\newcommand{\CheckShares}{\texttt{Check}}

\newcommand{\ASetup}{\texttt{ASetup}}
\newcommand{\AEval}{\texttt{AEval}}
\newcommand{\AWit}{\texttt{AWit}}
\newcommand{\AVer}{\texttt{AVer}}

\newcommand{\CSetup}{\texttt{CSetup}}
\newcommand{\CCommit}{\texttt{CCommit}}
\newcommand{\COpen}{\texttt{COpen}}

\newcommand{\SGen}{\texttt{SGen}}
\newcommand{\SSign}{\texttt{SSign}}
\newcommand{\SVer}{\texttt{SVer}}

\newcommand{\SoKSetup}{\texttt{SoKSetup}}
\newcommand{\SoKSign}{\texttt{SoKSign}}
\newcommand{\SoKVer}{\texttt{SoKVer}}

\newcommand{\TLEnc}{\texttt{TLEnc}}
\newcommand{\TLDec}{\texttt{TLDec}}

\newcommand{\frameindent}{\hspace*{\baselineskip}}

\colorlet{gris}{black!50!white}
\def\engris#1{\textcolor{gris}{#1}}

\colorlet{jaune}{yellow!90!black}
\def\enjaune#1{\textcolor{jaune!70!black}{#1}}

\colorlet{vert}{green!45!black}
\def\envert#1{\textcolor{vert}{#1}}

\colorlet{bleu}{blue!70!black}
\def\enbleu#1{\textcolor{bleu}{#1}}

\colorlet{rouge}{red!75!black}
\def\enrouge#1{\textcolor{rouge}{#1}}

\newcommand{\rouge}{\only{\color{rouge}}}

\xdefinecolor{violet}{rgb}{.5,0,.7}
\def\enviolet#1{\textcolor{violet}{#1}}

\xdefinecolor{violetft}{rgb}{0.199,0., 0.398}
\def\envioletft#1{\textcolor{violetft}{#1}}

\xdefinecolor{monbleu}{rgb}{0,.2,.54}
\newcommand{\monbleu}{\only{\color{monbleu}}}

\newcommand{\Ccal}{\mathcal{C}}
\newcommand{\Kcal}{\mathcal{K}}
\newcommand{\Mcal}{\mathcal{M}}
\newcommand{\Tcal}{\mathcal{T}}
\newcommand{\Xcal}{\mathcal{X}}

\newcommand{\Aenc}{E}
\newcommand{\Pair}[2]{\langle{#1},{#2}\rangle}
\newcommand{\Pk}[1]{\mathsf{pk}_{#1}}
\newcommand{\Sk}[1]{\mathsf{sk}_{#1}}
\newcommand{\h}{H}

\usepackage{wasysym}
\usepackage{amsmath}
\usepackage{amssymb}
\usepackage{amsfonts}
\usepackage{ragged2e}
%\usepackage{msc}
\usepackage{color}
\usepackage{stmaryrd}
\usepackage{eurosym}



\title{Cryptographic protocols}

\author{Markulf Kohlweiss\\
  School of Informatics \\
  University of Edinburgh}

\date{}

\begin{document}

\begin{frame}
  \maketitle

\end{frame}

\begin{frame}
  \frametitle{Context}
  Applications exchanging \enbleu{sensitive data} over a \enbleu{public network}:
  \begin{itemize}
  \item eBanking,
  \item eCommerce,
  \item eVoting,
  \item ePassports,
  \item cloud everything,
  \item \dots  \hfill cryptography is everywhere!
  \end{itemize}
  \pause
  \vspace{0.3cm}

  A malicious agent can:
  \begin{itemize}
  \item record, alter, delete, insert, redirect, reorder, and reuse past or current messages, and inject new messages 

    \enrouge{$\longrightarrow$ the network is the attacker}
  \item control dishonest participants
  \end{itemize}
  
\end{frame}

\begin{frame}
  \frametitle{The attacker controls the network (1)}
\vspace{-0.5cm}
  \begin{center}
    \includegraphics[scale=0.48]{images/TraceRoute1.pdf}
  \end{center}

\end{frame}

\begin{frame}
  \frametitle{The attacker controls the network (2)}
  \vspace{-0.5cm}

  \begin{columns}
    \begin{column}{0.4\textwidth}
      \vspace{-1cm}
       \begin{itemize}
       \item
       In July 2013, the NSA was accused of wiretapping large parts of data on the German Internet Exchange Point DE-CIX which was denied by Level 3. 
       
       \item A few months later, Level 3 was accused of tapping connections between Google and Yahoo data centers.
       \end{itemize}

       \url{https://en.wikipedia.org/wiki/Level_3_Communications}
    \end{column}

    \begin{column}{0.6\textwidth}
  \begin{center}
    \includegraphics[scale=0.30]{images/SnowdenRevelations.pdf}
  \end{center}
\end{column}
\end{columns}

\end{frame}

\begin{frame}
  \frametitle{The attacker controls the network (3)}
  \vspace{-0.5cm}

  \begin{columns}
    \begin{column}{0.4\textwidth}
      \vspace{-1cm}
       \begin{itemize}
       \item
       In July 2013, the NSA was accused of wiretapping large parts of data on the German Internet Exchange Point DE-CIX which was denied by Level 3. 
       
       \item A few months later, Level 3 was accused of tapping connections between Google and Yahoo data centers.
       \end{itemize}

       \url{https://en.wikipedia.org/wiki/Level_3_Communications}
    \end{column}

    \begin{column}{0.6\textwidth}
  \begin{center}
    \includegraphics[scale=0.4]{images/TraceRoute2.pdf}
  \end{center}
\end{column}

\end{columns}
\end{frame}


\iffalse
\begin{frame}
  \frametitle{All messages can be intercepted by an attacker (1)}

    \begin{center}
    \includegraphics[scale=0.85]{images/LoginForm.pdf}
  \end{center}
\end{frame}

\begin{frame}
  \frametitle{All messages can be intercepted by an attacker (2)}

  \begin{center}
    \includegraphics[scale=0.5]{images/WiresharkPasswordCapture.pdf}
  \end{center}
  \pause

  \begin{block}{}
    \enrouge{Man in the Middle (MitM) -} An attacker can \enrouge{intercept} packets, but also \enrouge{alter}, \enrouge{forge new}, and \enrouge{inject} packets
  \end{block}
\end{frame}
\fi

\begin{frame}
  \frametitle{More complex systems needed...}

  \onslide<2->{  
    \[
    \begin{array}[c]{ccc}
      \begin{array}[c]{c}
        \hspace{-0.75cm}\includegraphics<1->[scale=0.25]{images/su-mei-tse.pdf}
      \end{array}
      &
      \begin{array}[c]{c}
        \hspace{-0.5cm}\xrightarrow[m=MAC(K_M, E(K_E, \text{Transfer 100 } \footnotesize{\geneuro{}} \text{ on Amazon's account}))]{e=E(K_E, \text{Transfer 100 } \footnotesize{\geneuro{}} \text{ on Amazon's account})}
      \end{array}
      &
      \begin{array}[c]{c}
        \hspace{-0.5cm}\includegraphics<1->[scale=0.4]{images/RBS-logo.pdf}
      \end{array}
    \end{array}
    \]}

  \onslide<3->{
    \enbleu{Replay attack}
  
    \[
    \begin{array}[c]{ccccc}
      \hspace{0.2cm}\begin{array}[c]{c}
        \phantom{\includegraphics<1-2>[scale=0.25]{images/su-mei-tse.pdf}}
        \includegraphics<3->[scale=0.25]{images/su-mei-tse.pdf}
      \end{array}
      & 
      \hspace{0.2cm}\begin{array}[c]{c}
        \xrightarrow[]{(e, m)}
      \end{array}
      &
      \hspace{0.2cm}\begin{array}[c]{c}
        \phantom{\includegraphics<1-2>[scale=0.25]{images/intruder.pdf}}
        \includegraphics<3->[scale=0.25]{images/intruder.pdf}
      \end{array}
      &
      \hspace{0.2cm}\begin{array}[c]{c}
        \xrightarrow[]{(e, m)}
        \\
        \vdots \\
        \xrightarrow[]{(e, m)}
      \end{array}
      &
      \hspace{0.2cm}\begin{array}[c]{c}
        \phantom{\includegraphics<1-2>[scale=0.4]{images/RBS-logo.pdf}}
        \includegraphics<3->[scale=0.4]{images/RBS-logo.pdf}
      \end{array}
    \end{array}
    \]}

\end{frame}

\begin{frame}
  \frametitle{... to achieve more complex properties}
    \vspace{0.2cm}
    
  \begin{itemize}

  \item \enbleu{Confidentiality:} Some information should never be revealed to unauthorised entities.
    \vspace{0.2cm}

  \item \enbleu{Integrity}: Data should not be altered in an unauthorised manner since the time it was created, transmitted or stored by an authorised source.
    \vspace{0.2cm}

  \item \enbleu{Authentication:} Ability to know with certainty the identity of an communicating entity.
    \vspace{0.2cm}

  \item \enbleu{Anonymity:} The identity of the author of an action (\emph{e.g.} sending a message) should not be revealed.
    \vspace{0.2cm}

  \item \enbleu{Unlinkability:} An attacker should not be able to deduce whether different services are delivered to the same user
    \vspace{0.2cm}

  \item \enbleu{Non-repudiation:} The author of an action should not be able to deny having triggered this action.
    \vspace{0.2cm}

  \item \dots

  \end{itemize}
\end{frame}

\begin{frame}

  \frametitle{Cryptographic protocols}

  \begin{block}{Cryptographic protocols}
    A distributed algorithm employing \enbleu{cryptographic primitives} \\ \hfill and whose goal is \enbleu{``secure'' communication (or computation)}.
  \end{block}
  
  \vspace{0.5cm}
  \pause

  
  \begin{alertblock}{But!}
    Many exploitable errors are due not to design errors in the primitives, but to the way they are used, \emph{i.e.} bad protocol design and buggy or not careful enough implementation
  \end{alertblock}
  
\end{frame}

\begin{frame}
  \frametitle{Numerous deployed protocols are flawed...}
  \vspace{-0.5cm}

  \hfill ... and end up in the news :(
  \begin{center}
    \includegraphics[scale=0.40]{images/news.pdf}
  \end{center}

\end{frame}

\begin{frame}
  \frametitle{Logical attacks}
  \begin{center}
    \enrouge{Many of these attacks do not even break the cryptographic primitives!!}
        \includegraphics[scale=0.40]{images/backdoor.png}
  \end{center}
\end{frame}

\iffalse

\begin{frame}
  \frametitle{Example of a logical attack}

  \vspace{-0.5cm}

  Assume a commutative symmetric encryption scheme 
  \[
  \{\{m\}_{\enrouge{k_1}}\}_{\enrouge{k_2}} = \{\{m\}_{\enrouge{k_2}}\}_{\enrouge{k_1}}
  \]
  where $\{m\}_k$ denotes the encryption of message $m$ under the key $k$

  \envert{Example: stream ciphers}

  \begin{center}
    \footnotesize{
    \begin{tabular}[t]{ccccccccccc}
      & & & & & & \multicolumn{3}{c}{\onslide<6->{\enrouge{No authentication!}}} \\
      \onslide<2->{$A$} & & \onslide<2->{$B$} & & \onslide<6->{\enrouge{\vline}} & & \onslide<7->{$A$} & & \onslide<7->{\enrouge{$I$}} \\
      \onslide<2->{\vline} & & \onslide<2->{\vline} & & \onslide<6->{\enrouge{\vline}} & & \onslide<7->{\vline} & & \onslide<7->{\vline} \\
       \onslide<3->{\tiny{$\mathsf{new}\ K_A$}}  & & \onslide<2->{\vline} & & \onslide<6->{\enrouge{\vline}} & & \onslide<8->{\tiny{$\mathsf{new}\ K_A$}} & & \onslide<7->{\vline} \\
      \onslide<2->{\vline} & \onslide<3->{$\xrightarrow[\phantom{\{\{\text{pin: 3443}\}_{K_A}\}_{K_B}}]{\{\text{pin:3443}\}_{K_A}}$} & \onslide<2->{\vline} & & \onslide<6->{\enrouge{\vline}} & & \onslide<7->{\vline} & \onslide<8->{$\xrightarrow[\phantom{\{\{\text{pin: 3443}\}_{K_A}\}_{K_B}}]{\{\text{pin: 3443}\}_{K_A}}$} & \onslide<7->{\vline}\\
      \onslide<2->{\vline} & & \onslide<3->{\tiny{$\mathsf{new}\ K_B$}} & & \onslide<6->{\enrouge{\vline}} & & \onslide<7->{\vline} & & \onslide<9->{\tiny{$\mathsf{new}\ K_I$}} \\
      \onslide<2->{\vline} & \onslide<4->{$\xleftarrow[\phantom{\{\{\text{pin: 3443}\}_{K_A}\}_{K_B}}]{\{\{\text{pin: 3443}\}_{K_A}\}_{K_B}}$} & \onslide<2->{\vline} & & \onslide<6->{\enrouge{\vline}} & & \onslide<7->{\vline} & \onslide<9->{$\xleftarrow[\phantom{\{\{\text{pin: 3443}\}_{K_A}\}_{K_B}}]{\{\{\text{pin: 3443}\}_{K_A}\}_{K_I}}$} & \onslide<7->{\vline}\\
      \onslide<2->{\vline} & \onslide<5->{$\xrightarrow[\phantom{\{\{\text{pin: 3443}\}_{K_A}\}_{K_B}}]{\{\text{pin: 3443}\}_{K_B}}$} & \onslide<2->{\vline} & & \onslide<6->{\enrouge{\vline}} & & \onslide<7->{\vline} & \onslide<10->{$\xrightarrow[\phantom{\{\{\text{pin: 3443}\}_{K_A}\}_{K_B}}]{\{\text{pin: 3443}\}_{K_I}}$} & \onslide<7->{\vline}\\
    \end{tabular}}
  \end{center}
  \onslide<5->{\mbox{since $\{\{\text{pin: 3443}\}_{K_A}\}_{K_{B}} = \{\{\text{pin: 3443}\}_{K_{B}}\}_{K_{A}}$ by commutativity}}
  
\end{frame}

\fi

\title{Authentication and key agreement protocols}
\covertopicname{Authentication and key agreement protocols}
\author{}
\date{}

\begin{frame}
  \maketitle
\end{frame}

\begin{frame}
  \frametitle{Authentication and key agreement}

  \begin{itemize}
  \item Long-term keys should be used as little as possible to reduce ``attack-surface''
    \medskip{}
  \item The use of a key should be restricted to a specific purpose

    \envert{e.g. you shouldn't use the same RSA key both for encryption and signing}
    \medskip{}

  \item Public key algorithms tend to be computationally more expensive than symmetric key algorithms
    \medskip{}
  \item[\enrouge{$\rightsquigarrow$}] Long-term keys are used to establish short-term \enrouge{session keys}

    \envert{\emph{e.g.} TLS over HTTP, SSH for shell, Basic access control (BAC) protocol for ePassports, \emph{etc.}}
  \end{itemize}
\end{frame}

\begin{frame}
  \frametitle{Needham-Schroeder Public Key (NSPK)}

  NSPK: authentication and key agreement protocol
  \medskip{}

  \begin{center}
    \begin{tabular}[t]{ccccccc}
     & $A$ & & $B$ \\
     & \vline & & \vline \\
      \onslide<2->{\footnotesize{$\mathsf{new}\ N_A$}} & \vline & & \vline \\
     & \vline & \onslide<3->{$\xrightarrow[\hspace{4cm}]{\Aenc(\Pk{B}, \Pair{N_A}{A})}$} & \vline \\
     & \vline & & \vline & \onslide<4->{\footnotesize{$\mathsf{new}\ N_B$}} \\
     & \vline & \onslide<5->{$\xleftarrow[\hspace{4cm}]{\Aenc(\Pk{A}, \Pair{N_A}{N_B})}$} & \vline \\
     & \vline & \onslide<6->{$\xrightarrow[\hspace{4cm}]{\Aenc(\Pk{B}, \Pair{N_B}{A})}$} & \vline \\
      \\
      & \onslide<7->{$^{k_{AB}\leftarrow \h(N_A, N_B)}$} & & \onslide<7->{$^{k_{AB}\leftarrow \h(N_A, N_B)}$}
    \end{tabular}
  \end{center}
  \footnotesize{\engris{Roger M. Needham and Michael D. Schroeder. 1978. Using encryption for authentication in large networks of computers. Commun. ACM 21, 12 (Dec. 1978)}}
\end{frame}

\begin{frame}
  \frametitle{NSPK: security requirements}

  \begin{itemize}
  \item \enbleu{Authentication:} if Alice has completed the protocol, apparently with Bob, then
    Bob must also have completed the protocol with Alice.
    \medskip{}
  \item \enbleu{Authentication:} If Bob has completed the  protocol, apparently with Alice, then Alice must have completed the protocol with Bob.  
    \medskip{}
  \item \enbleu{Confidentiality:} Messages sent encrypted with the agreed key ($k\leftarrow \h(N_A,N_B)$) remain secret.
  \end{itemize}

\end{frame}

\begin{frame}
  \frametitle{NSPK: Lowe's attack on authentication}
  \bigskip{}

  \enrouge{Attack found 17 years after the publication of the NS protocol!!}
  \pause
  \small{
  \begin{center}
    \begin{tabular}[t]{ccccccc}
      & $A$ & & $I$ & & $B$ \\
      & \vline & & \vline & & \vline \\
      \onslide<2->{\footnotesize{$\mathsf{new}\ N_A$}} &  \vline  & & \vline & & \vline \\
      & \vline & \onslide<3->{$\xrightarrow[]{\Aenc(\Pk{I}, \Pair{N_A}{A})}$} & \vline & & \vline \\
      & \vline & & \vline & \onslide<4->{$\xrightarrow[]{\Aenc(\Pk{B}, \Pair{N_A}{A})}$} & \vline \\
      & \vline & & \vline & & \vline & \onslide<5->{\footnotesize{$\mathsf{new}\ N_B$}} \\
      & \vline & & \vline & \onslide<6->{$\xleftarrow[]{\Aenc(\Pk{A}, \Pair{N_A}{N_B})}$} & \vline \\
      & \vline & \onslide<7->{$\xleftarrow[]{\Aenc(\Pk{A}, \Pair{N_A}{N_B})}$} & \vline & & \vline \\
      & \vline & \onslide<8->{$\xrightarrow[\phantom{\Aenc(\Pk{A}, \Pair{N_A}{N_B})}]{\Aenc(\Pk{I}, \Pair{N_B}{A})}$} & \vline & & \vline \\
      & \vline & & \vline & \onslide<9->{$\xrightarrow[]{\Aenc(\Pk{B}, \Pair{N_B}{A})}$} & \vline \\
      &\onslide<10->{\enrouge{$^{k_{AI}\leftarrow \h(N_A, N_B)}$}} & & \onslide<10->{\enrouge{$^{k_{AI}\leftarrow \h(N_A, N_B)}$}} & &  \\
      & & & \onslide<10->{\enrouge{$^{k_{AB}\leftarrow \h(N_A, N_B)}$}}& & \hspace{-1cm}\onslide<10->{\enrouge{$^{k_{AB}\leftarrow \h(N_A, N_B)}$}}
    \end{tabular}
  \end{center}}
\footnotesize{\engris{[G. Lowe. ``An attack on the Needham-Schroeder public key authentication protocol''. Information Processing Letters (November 1995)]}}
\end{frame}

\begin{frame}{NSPK: Lowe's attack on authentication}{An awkward conversation}
  \small{
  \begin{center}
    \begin{tabular}[t]{ccccccc}
      &\onslide<1->{\enrouge{$^{k_{AI}\leftarrow \h(N_A, N_B)}$}} & & \onslide<1->{\enrouge{$^{k_{AI}\leftarrow \h(N_A, N_B)}$}} & &  \\
      & & & \onslide<1->{\enrouge{$^{k_{AB}\leftarrow \h(N_A, N_B)}$}}& & \hspace{-1cm}\onslide<1->{\enrouge{$^{k_{AB}\leftarrow \h(N_A, N_B)}$}}\\ \smallskip
      & $A$ & & $I$ & & $B$ \\
      & \vline & & \vline & & \vline \\
%      \onslide<2->{\footnotesize{$\mathsf{new}\ N_A$}} &  \vline  & & \vline & & \vline \\
%      & \vline & \onslide<3->{$\xrightarrow[]{\Aenc(\Pk{I}, \Pair{N_A}{A})}$} & \vline & & \vline \\
%      & \vline & & \vline & \onslide<4->{$\xrightarrow[]{\Aenc(\Pk{B}, \Pair{N_A}{A})}$} & \vline \\
%      & \vline & & \vline & & \vline & \onslide<5->{\footnotesize{$\mathsf{new}\ N_B$}} \\
      & \vline & & \vline & \onslide<2->{$\xleftarrow[]{\Aenc(k_{AB}, \text{"I love you!"})}$} & \vline \\
      & \vline & \onslide<3->{$\xleftarrow[]{\Aenc(k_{AB}, \text{"I love you!"})}$} & \vline & & \vline \\
      & \vline & \onslide<4->{$\xrightarrow[]{\Aenc(k_{AI}, \text{"I really don't!"})}$} & \vline & & \vline \\
      & \vline & & \vline & \onslide<5->{$\xrightarrow[]{\Aenc(k_{AI}, \text{"I really don't!"})}$} & \vline \\
         \end{tabular}
  \end{center}}
  Alice dislikes Ivan, but likes Bob.\smallskip \\
  Bob thinks he is talking to Alice, but is actually talking to the intruder Ivan relaying messages to Alice. 
%\footnotesize{\engris{G. Lowe. ``An attack on the Needham-Schroeder public key authentication protocol''. Information Processing Letters (November 1995)}}
\end{frame}

\begin{frame}
  \frametitle{NSPK: Lowe's fix}

  \begin{block}{The Needham-Schroeder-Lowe (NSL) protocol}
    \begin{center}
      \begin{tabular}[t]{ccccccc}
        & $A$ & & $B$ \\
        & \vline & & \vline \\
        \footnotesize{$\mathsf{new}\ N_A$} & \vline & & \vline \\
        & \vline & $\xrightarrow[\hspace{4cm}]{\Aenc(\Pk{B}, \Pair{N_A}{A})}$ & \vline \\
        & \vline & & \vline & \footnotesize{$\mathsf{new}\ N_B$} \\
        & \vline & $\xleftarrow[\hspace{4cm}]{\Aenc(\Pk{A}, \Pair{N_A}{\Pair{N_B}{\enrouge{B}}})}$ & \vline \\
        & \vline & $\xrightarrow[\hspace{4cm}]{\Aenc(\Pk{B}, \Pair{N_B}{A})}$ & \vline \\
        \\
        & $^{k_{AB}\leftarrow \h(N_A, N_B)}$ & & $^{k_{AB}\leftarrow \h(N_A, N_B)}$
      \end{tabular}
    \end{center}
  \end{block}

\end{frame}

\begin{frame}
  \frametitle{Forward secrecy}

  \begin{itemize}
  \item The NSL protocol is secure against an attacker that controls the network. 
  \item What if Alice's and Bob's private keys get compromised?
  \item What if the government forces Alice and Bob to reveal their private keys?
  \item Can we still protect confidentiality? 
  \end{itemize}
  \bigskip{}

  \begin{block}{Forward secrecy}
    A protocol ensures \enrouge{forward secrecy}, if even if long-term keys are compromised, past sessions of the protocol are still kept confidential, and this even if an attacker actively interferred.   
  \end{block}
\end{frame}

\begin{frame}
  \frametitle{The Station-to-Station (StS) protocol}
  \vspace{-0.5cm}

      \begin{center}
      \begin{tabular}[t]{ccccccc}
        & $A$ & & $B$ \\
        & \vline & & \vline \\
        \onslide<2->{\footnotesize{$\mathsf{pick}\ p$}} & \vline & & \vline \\
        \onslide<3->{\footnotesize{$x\gets\mathbb Z_p$}} & \vline & & \vline \\
        & \vline & \onslide<4->{$\xrightarrow[\hspace{4cm}]{g,\ p,\ g^x}$} & \vline \\
        & \vline & & \vline & \onslide<5->{\footnotesize{$y\gets\mathbb Z_p$}} \\
        & \vline & \onslide<6->{$\xleftarrow[\hspace{4cm}]{g^y, \envert{Cert_B}, \enbleu{E(\h(g^{yx}),} \envert{S(\mathsf{sk}_B}, (g^y, g^x)\envert{)}\enbleu{)}}$} & \vline \\
        & \vline & \onslide<7->{$\xrightarrow[\hspace{4cm}]{\envert{Cert_A}, \enbleu{E(\h(g^{yx}),} \envert{S(\mathsf{sk}_A}, (g^x, g^y)\envert{)}\enbleu{)}}$} & \vline \\
      \end{tabular}
    \end{center}
    \onslide<2->{\begin{itemize}
    \item where $p$ is a large prime
    \item and $g$ a generator of $\mathbb Z_p^*$
    \end{itemize}}
  \onslide<8->{
  \begin{block}{}
    The StS ensures mutual authentication, key agreement, and forward secrecy
  \end{block}}
  
\end{frame}

\begin{frame}{In defense of Eve}
  \vspace{-0.5cm}
  \begin{center}
\includegraphics[scale=0.4]{images/eve.png}
  \end{center}
    \footnotesize{\engris{XKCD:} \url{https://xkcd.com/177/}} \hfill \engris{Alice and Bob:} \url{https://cryptocouple.com/}
\end{frame}

\end{document}
